\section{Implementation and reference contracts}
\label{app:implementation}
The implementation is pinned to a specific llama.cpp/GGML source and numerical runtime. A candidate's identity includes model and tokenizer bytes, source/build, runtime libraries, graph configuration, KV type, sampler, stop policy and input IDs. The reference is an execution contract, not a mathematical ideal. When a later experiment lacks a historical complete-logit reference, it must establish a prospective reference with its own identity rather than claim that the older experiment contained the missing evidence.

The phase guard applies to target prefill and continuation; the layout is part of graph-cache compatibility. Observation-disabled execution must use the same guard. The QB transformation does not replace native acceptance, correction, pending-logit consumption or KV trimming. Unsupported fused attention features and layouts remain outside the qualified domain.

\section{Original CUDA panel and numerical replay}
\label{app:cuda}
The original complete-output campaign covers 24 fixed inputs per family, with repeated observations retained as repeats rather than independent tasks. The implementation comparison uses a separate 12-input panel per family and seven repair/draft arms. The six-session E2 comparison uses fresh workers, counterbalanced PQ/QB blocks and 1,620 requests including 324 warmups. It preserves task-level regressions and does not use request-level bootstrap or outlier trimming.

The 72-context replay spans three model families, three layer positions, target prefill and continuation, and q=1 through q=4. All PQ/QB pairs match bitwise at QK, softmax and AV. The 18 agreements between ordinary multi-query attention and either repaired path are the q=1 common-path controls; the remaining 54 contexts first diverge at QK. This is local correspondence on captured operands, not a bound on model logits or a universal reduction-order result.

\section{RTX mapping interventions}
\label{app:rtxmap}
The CUDA Graph intervention compares verified graph capture/instantiate/replay with the corresponding non-replayed mode. QB completed-graph ON/OFF ratios of 0.615--0.849 show that execution mode changes the mapping-cost landscape. This does not isolate launch overhead from all other graph-mode effects.

The dispatch intervention holds logical length fixed while changing physical representation and masks. Logical length 1023 padded to physical 1024 reaches a three-kernel route but fails the original reference. The qualified 1023-to-1025 route remains eight kernels and is about 1.5\% more expensive. The rejected route has no valid performance estimate under the reference contract.

The compact-QH transaction includes the common Q preparation, upload/completion, graph compute/synchronization and complete ordinary readback. K/V and mask are resident; target buffers are preallocated. Natural Q production and KV updates are excluded. In the graph-ON panel, compact-QH/QB ratios are 1.002216 (q=2) and 1.001598 (q=4); the tested branch did not establish a practical advantage.

\section{B570 request panel and software counterfactual}
\label{app:b570}
The scoped Vulkan GQA path uses Qwen3-8B Q8_0 target and Qwen3-0.6B Q8_0 draft. The synthetic operator panel fixes Hq/Hkv=32/8, d=128, F16 K/V, F32 Q/mask, capacity 1024 and physical key extents 256/512. The first candidate changed AV evaluation and failed byte equality at key512/q2; the next candidate retained the original noncontiguous AV path and passed the six registered synthetic coordinates. These records do not establish arbitrary GQA support.

The four-input formal panel reports the following within-panel estimates.

\begin{table}[h]
\centering\small
\caption{B570 request-panel cost results.}
\begin{tabular}{@{}lrr@{}}
\toprule
Input & QB/dual-PQ service & QB/target-only \\
\midrule
HumanEval/38 & 0.99464 & 2.5978 \\
HumanEval/65 & 0.99206 & 2.6520 \\
HumanEval/115 & 0.99728 & 2.6198 \\
HumanEval/127 & 0.99666 & 2.6097 \\
\midrule
Geometric mean & 0.995144 & 2.619763 \\
\bottomrule
\end{tabular}
\end{table}

The primary layout effect is a 0.486\% point reduction, below the registered 2\% screen. The full service ratio to target-only is 2.619763. The panel includes 64 measured requests and 16 explicit warmups over four independent process sessions; the full request timer includes fresh-context creation and teardown, prefill, generation, state maintenance and synchronization, while shared model loading is separately recorded. Host phase intervals are not GPU kernel-active time.

The draft CPU parameter intervention changes only the draft context thread parameter from eight to one. All 346 saved target/draft logits vectors remain byte-identical to the earlier qualified QB result, and all endpoints and cycle traces pass. The DRAFT1/dual-baseline service ratio is 1.003994, so the candidate is about 0.399\% slower. Source inspection found one-task embedding operations and worker-count clamping; the experiment does not prove that fewer workers actually ran.

\section{RTX platform closure}
\label{app:rtxclosure}
The RTX closure is a prospective service study because the original E0 audit did not recover a complete historical logits reference for every selected input. The new driver/reference identity was frozen before the first model process, and qualification was sealed before clean timing. The model/tokenizer bytes, output cap and 21 numerical DLLs remained unchanged. This does not retrofit the missing reference into the historical PR104 study.

The fixed panel consists of four previously exposed inputs: MBPP/413 (34 tokens, EOG), GSM8K/test/120 (128 tokens, cap), GSM8K/test/322 (128 tokens, cap) and GSM8K/test/443 (128 tokens, cap). The primary study uses k=1, actual caller q=2, QB physical q=2, PQ serial q=1 and GQA32/8 across 36 target layers. All endpoints, decision-relevant logits and cycle/state traces pass.

\begin{table}[h]
\centering\small
\caption{RTX closure within-panel service ratios.}
\begin{tabular}{@{}lrr@{}}
\toprule
Input & QB/dual-PQ & QB/target-only \\
\midrule
MBPP/413 & 0.988025 & 0.608126 \\
GSM8K/120 & 1.004279 & 0.578678 \\
GSM8K/322 & 0.963695 & 0.610027 \\
GSM8K/443 & 1.001663 & 0.535651 \\
\midrule
Geometric mean & 0.989284 & 0.582325 \\
\bottomrule
\end{tabular}
\end{table}

The experiment comprises four independent timing processes, 64 measured requests and 64 separately charged warmups. The primary service timer includes fresh target/draft context and graph reservation, prefill, generation, greedy decisions, output/state handling, synchronization and teardown. Both model weights remain resident for target-only R; model loading is separate. R retains target ubatch 1, while speculative target verification uses ubatch 4 and the draft ubatch 1.

The 1.072\% QB-over-dual-PQ point reduction does not clear the registered screen because the largest input-level A/A imbalance is 5.208\%, and two process-session aggregates favor PQ. The study therefore establishes no practical layout benefit and does not prove equivalence. In contrast, QB/R=0.582325 indicates 41.8\% lower service time than target-only for this finite panel and execution policy. It is not a production-serving ceiling or a universal speculative speedup.

\section{Cross-platform claim boundaries}
\label{app:boundaries}
The RTX and B570 panels use different frontends, devices, runtime contracts, reference identities, context policies and cost windows. Their service ratios answer separate within-panel questions. They must not be pooled into one mean, used to rank devices, or interpreted as evidence that a particular physical resource is the bottleneck.

The new RTX raw archive contains 2,900 raw vectors across 3,110 archive members and is 2,270,281,687 bytes locally. Its members were reopened and hash-checked on the originating host. The full raw archive was not externally authenticated in this review. GitHub stores compact records, hashes, receipts and estimator outputs; this supports compact-data reconstruction but is not independent re-reading of the full raw vectors.
